\documentclass[12pt,a4paper,twocolumn]{article}
\usepackage[utf8]{inputenc}
\usepackage[spanish]{babel}
\usepackage{amsmath}
\usepackage{amsfonts}
\usepackage{amssymb}
\usepackage{graphicx}
\usepackage{geometry}
\usepackage{hyperref}
\usepackage{booktabs}
\usepackage{caption}
\usepackage{float}
\usepackage{abstract}
\usepackage{titlesec}
\usepackage{authblk}

\geometry{top=2cm, bottom=2cm, left=1.5cm, right=1.5cm}

\titleformat{\section}{\large\bfseries}{\thesection.}{1em}{}
\titleformat{\subsection}{\normalsize\bfseries}{\thesubsection.}{1em}{}

\title{\textbf{Desarrollo de un Gemelo Digital Basado en Modelado Fenomenológico, para la Producción de Proteína Aislada de Soya (ISP)}}

\author[1]{Samuel Aguilera Araujo}
\affil[1]{Ingeniería Industrial}

\date{\today}

\begin{document}

\twocolumn[
  \begin{@twocolumnfalse}
    \maketitle
    \begin{abstract}
      Se desarrolló el Gemelo Digital (\textit{Digital Twin}) como pilar de la transformación industrial hacia la Industria 5.0, presentando su arquitectura estandarizada bajo ISO 23247. Sobre esta fundamentación, se detalla el diseño de un Gemelo Digital fenomenológico aplicado a la producción de Proteína Aislada de Soya (ISP) con capacidad de 1000 kg/h. El modelo acopla ecuaciones fenomenológicas no lineales de transferencia de masa y energía con tecnología avanzada de separación por membranas (UF/DF/RO) y un motor de restricciones operativas en tiempo real. Los resultados validan que la integración de una etapa híbrida de preconcentración por Ósmosis Inversa (OI) reduce el OPEX térmico del evaporador de doble efecto en un 25.4\% al retirar 2718.3 kg/h de agua en frío, requiriendo un consumo eléctrico de tan solo 9.0 kW. 
      
      \vspace{1em}
      \\ \textbf{Palabras clave:} Gemelo Digital, Ingeniería Industrial, PINNs, MPC, Proteína de Soya, Ósmosis Inversa, ISO 23247, FMEA Dinámico.
      \vspace{2em}
    \end{abstract}
  \end{@twocolumnfalse}
]

\section{Introducción}

La transición hacia la Industria 4.0 y 5.0 demanda abandonar los modelos heurísticos estáticos y adoptar infraestructuras ciberfísicas capaces de reducir la varianza operativa de forma autónoma. En el marco metodológico de Lean Six Sigma (LSS), el objetivo central es minimizar los defectos, optimizar la capacidad del proceso ($C_{pk}$) y garantizar que las Características Críticas para la Calidad (CTQ) se mantengan dentro de límites de especificación rigurosos, independientemente de las perturbaciones externas.

El Gemelo Digital (\textit{Digital Twin} -- DT) representa una evolución de gran impacto en la ingeniería industrial, consolidándose como herramienta clave para esta transición. A diferencia de los métodos clásicos de simulación estática, un gemelo digital proporciona una representación virtual dinámica y fiel de un sistema físico, permitiendo modelar su comportamiento ante perturbaciones y optimizar variables operativas en tiempo real.

\subsection{Clasificación de Modelos Digitales}
En el ámbito académico e industrial (Kritzinger et al., 2018), es fundamental diferenciar los niveles de integración virtual:
\begin{enumerate}
    \item \textbf{Modelo Digital (\textit{Digital Model}):} Representación virtual estática sin intercambio de datos automatizado. La transferencia es manual; el estado virtual no influye ni se actualiza automáticamente con el activo físico.
    \item \textbf{Sombra Digital (\textit{Digital Shadow}):} Flujo unidireccional y automatizado desde el activo físico hacia la réplica virtual. Cualquier cambio en la planta se refleja inmediatamente, pero el modelo no ejerce acciones de control automático.
    \item \textbf{Gemelo Digital (\textit{Digital Twin}):} Flujo bidireccional, continuo y automatizado. La réplica virtual recibe telemetría en tiempo real y ejecuta algoritmos de simulación, optimización y control predictivo, enviando comandos de vuelta al sistema físico para autorregular su desempeño operacional.
\end{enumerate}

\subsection{Funciones Operativas del Gemelo Digital}
Un gemelo digital actúa como un motor cognitivo:
\begin{itemize}
    \item \textbf{Monitoreo y Diagnóstico Continuo:} Integra datos de sensores IoT para diagnosticar la salud de equipos y detectar desviaciones.
    \item \textbf{Simulación de Escenarios (\textit{What-If}):} Predice y evalúa el impacto de perturbaciones dinámicas sin poner en riesgo la planta real.
    \item \textbf{Optimización Predictiva:} Ajusta continuamente los puntos de consigna (\textit{Set Points}) para maximizar rendimientos, reducir mermas y conservar la calidad del producto.
\end{itemize}

\subsection{Estructura del Gemelo Digital Aplicado a ISP}

El proceso se simula a través de las siguientes etapas operativas:
\begin{enumerate}
    \item \textbf{Etapa 0 (Acondicionamiento y Molienda):} Reducción de harina desgrasada a $<75$ $\mu$m.
    \item \textbf{Etapa 1 (Extracción Alcalina):} Solubilización de proteínas a pH 8.75 y 55$^\circ$C.
    \item \textbf{Etapa 1.2 (Clarificación Centrífuga):} Separación a 1800 g para descartar la fibra insoluble (Okara).
    \item \textbf{Etapa 2 (Pasteurización HTST):} Inactivación de factores antinutricionales a 80$^\circ$C por 22 s.
    \item \textbf{Etapa Híbrida (UF/DF/RO):} Fraccionamiento por membranas y preconcentración por Ósmosis Inversa a 24 bar.
    \item \textbf{Etapa 3 (Evaporación Térmica):} Evaporación de doble efecto al vacío ($0.40$ bar abs, $75^\circ$C).
    \item \textbf{Etapa 4 (Precipitación Isoeléctrica):} Neutralización ácida con HCl hasta pH 4.5.
    \item \textbf{Etapa 5 (Secado por Atomización):} Secado en Spray Dryer a 190$^\circ$C de entrada.
\end{enumerate}

\section{Marco Arquitectónico y Estandarización}

\subsection{Arquitectura ISO 23247}
La escalabilidad comercial y la interoperabilidad de un ecosistema ciberfísico exigen apego a normativas internacionales. La serie ISO 23247 (\textit{Digital Twin Framework for Manufacturing}) segmenta la arquitectura en cuatro dominios funcionales, garantizando trazabilidad absoluta y baja latencia computacional:

\begin{enumerate}
    \item \textbf{Capa de Elementos de Manufactura Observables (OME):} Representa los activos físicos: tanques agitados TK-101, bombas centrífugas sanitarias, intercambiadores de calor HX-201, centrífugas decantadoras, módulos de filtración por membranas y el secador por atomización SD-501. Los equipos deben cumplir con normativas ASME BPE y directrices EHEDG.
    \item \textbf{Entidad de Comunicación de Dispositivos:} Puente neuro-digital que recopila señales analógicas/digitales de los elementos OME y despacha programas de control hacia los actuadores. Su robustez depende de protocolos de telemetría como OPC UA, MQTT y MTConnect.
    \item \textbf{Entidad del Gemelo Digital:} Núcleo algorítmico y matemático donde residen los simuladores fenomenológicos y motores de validación de restricciones.
    \item \textbf{Entidad de Usuario:} Interfaces gráficas y consolas interactivas (e.g., Python/Streamlit, realidad aumentada) para la interacción con métricas e indicadores de rendimiento.
\end{enumerate}

\begin{table}[H]
\centering
\caption{Mapeo ISO 23247 a la Planta ISP}
\begin{tabular}{lll}
\toprule
Capa ISO & Componentes ISP & Protocolos \\
\midrule
OME & TK-101, P-102, SD-501 & ASME BPE \\
Comunicación & PLCs, Sensores IoT & OPC UA, MQTT \\
Gemelo Digital & PINNs, MPC, EKF & AAS (IEC 63278) \\
Usuario & Dashboard Streamlit & REST, WebSocket \\
\bottomrule
\end{tabular}
\end{table}

\subsection{Asset Administration Shell (AAS)}
El estándar AAS (IEC 63278) actúa como la envoltura digital que estandariza propiedades, parámetros y servicios de cada activo físico. El metamodelo organiza la información mediante ``submodelos'' (\textit{Submodels}), referenciando vocabularios como ECLASS e IEC CDD. La implementación del AAS se integra mediante OPC UA a través del mapeo I4AAS, proveyendo interpretación semántica de las variables de proceso. Cuando el AAS se publica sobre un servidor OPC UA, el algoritmo del Gemelo Digital no lee simplemente un ``Tag\_001 = 4.5''; interpreta semánticamente que la variable corresponde al ``pH del tanque de precipitación isoeléctrica TK-401'', evaluando su margen contra las restricciones modeladas.

\section{Metodología de Desarrollo}

\subsection{Arquitectura del Software en Tres Capas}
El software del Gemelo Digital se diseñó con estructura modular en Python:
\begin{enumerate}
    \item \textbf{Capa de Fenomenología (Core):} Resuelve los sistemas de ecuaciones no lineales que describen las operaciones unitarias. Utiliza correlaciones termodinámicas y cinéticas de transferencia de masa y energía (archivo \texttt{core/stage\_equations.py}).
    \item \textbf{Capa de Restricciones (\textit{Constraints Engine}):} Evalúa en cada paso de tiempo si las variables de operación violan capacidades físicas o mecánicas de los equipos (archivo \texttt{core/equipment\_constraints.py}).
    \item \textbf{Capa de Visualización y Control (Streamlit):} Diseñada para la sala de control digital interactiva (\texttt{app.py}), mostrando KPIs en tiempo real, gráficos dinámicos y matrices FMEA.
\end{enumerate}

\subsection{Modelo de Inercia Temporal y Retardo}
Para simular el comportamiento transitorio de una planta física real, el Gemelo Digital implementa un modelo de retardo de primer orden (inercia térmica y de mezcla) para calcular las variables de proceso reales (\textit{Process Variables} -- PV) en cada paso de tiempo $\Delta t$:
\begin{equation}
    PV(t + \Delta t) = PV(t) + [SP - PV(t)] \cdot (1 - e^{-\Delta t / \tau})
\end{equation}
Donde $\tau$ es la constante de tiempo característica de la variable. Para variables térmicas y de volumen, se fija nominalmente en $\tau = 20.0$ segundos.

\section{Modelado Fenomenológico por Etapas}

El motor matemático del Gemelo Digital resuelve secuencialmente los balances para las 6 fases del proceso industrial (Etapa 0 a Etapa 5).

\subsection{Etapa 0: Acondicionamiento y Molienda}
Se asume una alimentación de harina de soya desgrasada $\dot{m}_{\text{soya}} = 1000$ kg/h con fracción proteica $X_p = 0.375$. La masa de proteína entrante es:
\begin{equation}
    \dot{m}_{p,\text{in}} = \dot{m}_{\text{soya}} \cdot X_p = 375.0 \text{ kg/h}
\end{equation}
Se fija una relación solvente/sólido de 12:1, resultando en un caudal de agua de $\dot{m}_{\text{agua}} = 12000$ kg/h ($Q = 12.0$ m$^3$/h). La potencia hidráulica teórica de la bomba P-101 se modela como:
\begin{equation}
    P_{\text{bomba}} = \frac{\rho \cdot g \cdot Q \cdot H}{3.6 \times 10^6 \cdot \eta_{\text{tot}}}
\end{equation}
Donde $\rho = 1000$ kg/m$^3$, $g = 9.81$ m/s$^2$, $H = 10.7$ m de altura manométrica y la eficiencia total es $\eta_{\text{tot}} = \eta_{\text{hid}} \cdot \eta_{\text{motor}} = 0.65 \cdot 0.90 = 0.585$. El molino de impacto reduce la harina a un tamaño de partícula de 200 mesh ($<75$ $\mu$m).

\subsection{Etapa 1: Extracción y Lixiviación Alcalina}
La lixiviación de proteínas globulares se realiza a pH 8.75 y $55^\circ$C. La eficiencia de extracción ($\eta_{\text{ext}}$) se modela mediante factores multiplicativos no lineales:
\begin{equation}
    \eta_{\text{ext}} = \eta_{\text{base}} \cdot f_{\text{pH}} \cdot f_T \cdot f_{\tau} \cdot f_R
\end{equation}
Donde $\eta_{\text{base}} = 0.88$. El factor de pH modela la solubilidad en ``U'':
\begin{equation}
    f_{\text{pH}} = 1.0 - 0.05 \cdot |\text{pH} - 8.75|
\end{equation}
El factor térmico $f_T$ es:
\begin{equation}
    f_T = 1.0 - 0.005 \cdot |T - 55.0|
\end{equation}
El factor de tiempo de residencia en el tanque agitado TK-101 es:
\begin{equation}
    f_{\tau} = \min\left(1.0, \frac{\tau_{\text{res}}}{60.0}\right)
\end{equation}
La inyección de NaOH comercial al 20\% p/p para control de pH añade $60.0$ kg/h de solución alcalina. El flujo másico total del lodo considerando pérdidas operacionales mecánicas ($F_{\text{loss}} = 0.02$) es:
\begin{equation}
    \dot{m}_{\text{lodo}} = (\dot{m}_{\text{soya}} + \dot{m}_{\text{agua}} + \dot{m}_{\text{NaOH}}) \cdot (1 - F_{\text{loss}}) = 12740 \text{ kg/h}
\end{equation}
El número de Reynolds del agitador en el TK-101 ($\text{Re}_a$) y la potencia consumida son:
\begin{equation}
    \text{Re}_a = \frac{\rho \cdot N \cdot D_a^2}{\mu} = \frac{1050 \cdot 1.33 \cdot 1.08^2}{0.020} = 81424
\end{equation}
\begin{equation}
    P_{\text{agitador}} = N_p \cdot \rho \cdot N^3 \cdot D_a^5 = 1.3 \cdot 1050 \cdot 1.33^3 \cdot 1.08^5 = 4.71 \text{ kW}
\end{equation}

El modelo mecanicista dictamina que si el pH cae a 8.0 en TK-101, la eficiencia se contrae desde el 88.0\% hasta un 75.9\%, colapsando la masa proteica soluble y deprimiendo la producción de polvo en 12.1\%. El MPC sincroniza dinámicamente la dosificación de NaOH con las lecturas de masa de la tolva.

\subsection{Etapa 1.2: Clarificación Centrífuga}
El lodo se separa en una centrífuga decantadora a 1800 g. Siguiendo la Ley de Stokes para campo centrífugo:
\begin{equation}
    v_s = \frac{d^2 (\rho_p - \rho_f) r \omega^2}{18 \mu}
\end{equation}
Se obtiene una salida de $1645$ kg/h de Okara húmedo (65\% humedad) y $11095$ kg/h de extracto clarificado. Con la merma por transporte ($F_{\text{loss}} = 0.02$), el caudal de extracto neto que ingresa a pasteurización es de $10873.1$ kg/h, con una masa proteica útil de:
\begin{equation}
    \dot{m}_{p,\text{extracto}} = 375.0 \cdot 0.88 \cdot 0.98^2 = 316.9 \text{ kg/h}
\end{equation}

\subsection{Etapa 2: Pasteurización HTST}
Se modela un intercambiador de placas (PHE) a $80^\circ$C. Con recuperación térmica del 55\% ($\text{HR} = 0.55$), el balance de energía determina la carga térmica de utilidad:
\begin{equation}
    \dot{Q}_{\text{utilidad}} = \dot{m}_{\text{extracto}} \cdot C_p \cdot (T_{\text{pasteur}} - T_{\text{in}}) \cdot (1 - \text{HR})
\end{equation}
Con $C_p = 3.9$ kJ/(kg$\cdot$K), $T_{\text{in}} = 25^\circ$C, $\dot{Q}_{\text{utilidad}} = 3.02 \text{ kg/s} \cdot 3.9 \cdot 55 \cdot 0.45 = 291.5$ kW. El área requerida es de $12.9$ m$^2$ ($U = 2000$ W/m$^2$K), utilizándose un equipo comercial de $39.6$ m$^2$ para absorber factores de ensuciamiento.

\subsection{Etapa Híbrida: Tecnología de Membranas (UF/DF/RO)}

\subsubsection{Ultrafiltración y Diafiltración (UF/DF)}
La inserción de tecnología avanzada de separación por membranas representa un cambio de paradigma hacia la circularidad y la eficiencia energética. El Gemelo Digital controla flujos cruzados a través de módulos de membranas con cortes de peso molecular (MWCO) desde 20{,}000 hasta 300{,}000 Daltons. Este fraccionamiento permite que el permeado arrastre agua, azúcares solubles, ácido fítico y oligosacáridos, reteniendo y concentrando las fracciones macro-proteínicas sin cambios de fase térmicos ni desnaturalización inducida por ácidos.

La adición de ciclos de Diafiltración continuos es crítica. Operar manteniendo sólidos retenidos cercanos al 10\% (w/w) maximiza la permeabilidad volumétrica ($J_w$). En términos de CTQs, las metodologías de membranas logran niveles de remoción del 86.5\% en carbohidratos y 64.5\% en cenizas (Mondor et al., 2009), resultando en un producto de pureza funcional $>$90\% (base seca) y una recuperación proteica entre un 17\% y un 26\% estructuralmente superior a la tecnología puramente isoeléctrica (Alibhai et al., 2006).

\subsubsection{Ósmosis Inversa (RO)}
La Ósmosis Inversa actúa como etapa de preconcentración. El transporte másico de solvente a través de la membrana de poliamida obedece al modelo simplificado de Kedem-Katchalsky (reflexión $\sigma \approx 1$):
\begin{equation}
    J_w = A_w \cdot (\text{TMP} - \Delta \pi)
\end{equation}
A una Presión Transmembrana (TMP) de 24 bar, la OI retira un caudal de permeado puro de $\dot{m}_{\text{permeado}} = 2718.3$ kg/h. La potencia de la bomba de alta presión es de:
\begin{equation}
    P_{\text{bomba, RO}} = \frac{Q_{\text{feed}} \cdot \text{TMP}}{\eta_{\text{pump}}} = \frac{0.00288 \cdot 2.4 \times 10^6}{0.80} = 9.0 \text{ kW}
\end{equation}
El ahorro de calor sensible y latente en el evaporador subsecuente es:
\begin{equation}
    \dot{Q}_{\text{saved}} = \dot{m}_{\text{permeado}} \cdot \frac{\lambda_{\text{vap}}}{3600} = 2718.3 \cdot \frac{2315}{3600} = 945.5 \text{ kW}
\end{equation}

\subsubsection{Modelado Predictivo de Ensuciamiento}
El impedimento técnico primordial en la filtración de biopolímeros es la polarización por concentración y el \textit{fouling}. Los modelos puramente analíticos (Ley de Darcy) subestiman los efectos de repulsión electrostática de solutos. Mediante redes neuronales convolucionales (CNN) o perceptrones multicapa informados por primeros principios (FFNN), el Gemelo Digital es entrenado con datos históricos para inferir la resistencia de la membrana. Cuando las arquitecturas LSTM proyectan saturación inminente de la microestructura poral, el controlador MPC modula la velocidad de flujo cruzado, ajusta la TMP o determina autónomamente el punto óptimo de transición entre ciclos de ultrafiltración y lavado químico (CIP), evitando pérdidas irreversibles de permeabilidad.

\subsection{Etapa 3: Evaporación Térmica de Doble Efecto}
El evaporador opera al vacío ($0.40$ bar abs, $T_{\text{ebullición}} \approx 75^\circ$C) para evitar la reacción de Maillard y conservar la calidad proteica. El balance másico se calcula para concentrar hasta $x_p = 0.23$ de sólidos:
\begin{equation}
    \dot{m}_{\text{concentrado}} = \frac{\dot{m}_{\text{solutos}}}{x_p} = \frac{304.4 \text{ kg/h}}{0.23} = 1323.5 \text{ kg/h}
\end{equation}
La cantidad de agua evaporada térmicamente es:
\begin{equation}
    \dot{m}_{\text{evaporada}} = \dot{m}_{\text{feed}} - \dot{m}_{\text{concentrado}} - \dot{m}_{\text{permeado, RO}}
\end{equation}
El consumo de vapor de caldera se calcula con la economía de doble efecto ($E = 1.85$):
\begin{equation}
    \dot{m}_{\text{vapor}} = \frac{\dot{m}_{\text{evaporada}}}{1.85}
\end{equation}

\subsection{Etapa 4: Precipitación Isoeléctrica}
El concentrado es acidificado con HCl al 10\% p/p ($110.0$ kg/h de solución) hasta alcanzar pH 4.5. En este punto la repulsión electrostática colapsa ($\zeta \approx 0$ mV) y la solubilidad es mínima. La eficiencia de precipitación es:
\begin{equation}
    \eta_{\text{precip}} = 0.98 \cdot f_{\text{pH,precip}}
\end{equation}
Si la desviación de pH supera el rango crítico ($4.5 \pm 0.4$), se penaliza drásticamente el rendimiento. El diámetro medio de los flóculos ($d$ en $\mu$m) se modela empíricamente:
\begin{equation}
    d = 320.0 - 250.0 \cdot |\text{pH} - 4.5| + 1.8 \cdot (\tau_{\text{precip}} - 25.0)
\end{equation}
La pasta de proteína con $50\%$ de humedad se separa en la centrífuga CF-401 con eficiencia de recuperación sólida del $98\%$.

\subsection{Etapa 5: Secado por Atomización (Spray Dryer)}

\subsubsection{Balance Entálpico}
La pasta se atomiza con aire seco a $190^\circ$C. El balance psicrométrico restringe la evaporación para evitar desnaturalización (temperatura de gota limitada por el bulbo húmedo a $<70^\circ$C). El polvo se retira a humedad residual de 5\%, obteniéndose una producción neta de $301.6$ kg/h de ISP con rendimiento proteico global del 76.4\%.

\subsubsection{Modelado CFD y Transición Vítrea ($T_{\text{g}}$)}
El Gemelo Digital trasciende las suposiciones homogéneas mediante sub-modelos de Dinámica de Fluidos Computacional (CFD) acoplados con Modelos de Balance Poblacional (PBM) y aproximaciones cinéticas especializadas (REA -- \textit{Reaction Engineering Approach}). Estas ecuaciones simulan la contracción de diámetro de las gotas, los gradientes térmicos y la evaporación de multicomponentes en un entorno compartimentalizado (1.5-D).

El cuello de botella crítico en el \textit{spray drying} de alimentos recae en la adherencia parasitaria del polvo a las paredes de la torre (\textit{wall deposition}). Este fenómeno es modulado por la Temperatura de Transición Vítrea ($T_{\text{g}}$) del polímero de soya. La física de polímeros (gobernada por las cinéticas WLF -- Williams-Landel-Ferry) establece que cuando un sólido amorfo sobrepasa su $T_{\text{g}}$, transita desde un estado cristalino-vítreo hacia un estado gomoso-plástico donde la viscosidad superficial colapsa. Se define como \textit{Sticky Point} aquel umbral térmico situado típicamente entre $10$ y $20$ K por encima de $T_{\text{g}}$ (Roos, 1995; Bhandari \& Howes, 1999).

Puesto que el agua libre actúa como plastificante, cualquier incremento en la humedad residual deprime brutalmente el $T_{\text{g}}$. Si el Gemelo Digital detecta, a través de perfiles CFD y sensores de humedad, que las partículas se aproximan al límite de $\Delta T = (T_{\text{partícula}} - T_{\text{g}}) > 15$ K, la algoritmia de IA predice colapso inminente de deposición adhesiva. El controlador MPC reduce la carga de alimentación o ajusta las tasas psicrométricas de inyección, manteniendo un régimen operacional seguro.

\section{Restricciones Operativas e Hidráulica}

\subsection{Motor de Verificación de Restricciones}
El Gemelo Digital ejecuta iterativamente el algoritmo de validación de capacidad de equipos. Las restricciones mecánicas y de proceso clave se detallan en la Tabla 4.

\begin{table}[H]
\centering
\caption{Límites de Capacidad y Restricciones del Proceso}
\begin{tabular}{lcc}
\toprule
Restricción / Equipo & Variable Crítica & Límite Máximo \\
\midrule
Tanque Agua TK-100 & Fracción Llenado & 90\% \\
Centrífuga CF-102 & Caudal Hidráulico & 16.0 m$^3$/h \\
Pasteurizador HX-201 & Carga Térmica & 100\% Cap. \\
Evaporador EV-301 & Capacidad Evap. & 10.0 m$^3$/h \\
Secador Spray SD-501 & Capacidad Agua & 450 kg/h \\
Secador SD-501 & Carga Específica & 1.5 kg/(h$\cdot$m$^3$) \\
\bottomrule
\end{tabular}
\end{table}

\subsection{Cavitación de Bombas y Margen de NPSH}
Para prevenir la cavitación de la bomba P-102 que opera con extracto caliente a $55^\circ$C, el Gemelo Digital calcula dinámicamente el NPSH disponible ($\text{NPSH}_a$):
\begin{equation}
    \text{NPSH}_a = \frac{P_{\text{atm}} - P_v(55^\circ\text{C})}{\rho \cdot g} + H_s - h_f
\end{equation}
Con $P_{\text{atm}} = 101{,}325$ Pa, la presión de vapor saturado del agua a $55^\circ$C es $P_v = 15{,}700$ Pa. Con altura estática de succión $H_s = 2.0$ m y pérdidas por fricción en tubería de succión $h_f = 0.7$ m:
\begin{equation}
    \text{NPSH}_a = \frac{101325 - 15700}{1050 \cdot 9.81} + 2.0 - 0.7 = 9.6 \text{ m}
\end{equation}
Dado que el fabricante especifica un $\text{NPSH}_r \approx 2.5$ m, la operación cuenta con un margen de seguridad de $7.1$ m, garantizando prevención de daño mecánico en álabes y desnaturalización por vibración sónica. El MPC evalúa en cada diferencial de tiempo este margen; si prevé que la aceleración de la bomba ocasione una caída de presión por debajo de la curva de ebullición, interviene reduciendo los hertz del variador, evadiendo la cavitación destructiva.

\subsection{Diseño Higiénico y Filosofía SMED}
La red de tuberías de acero AISI 316L está dimensionada para velocidades de flujo entre $0.8 - 1.2$ m/s, mitigando el cizallamiento proteico. Para la fase de limpieza CIP (\textit{Clean In Place}), las bombas sanitarias elevan la velocidad a $v \ge 1.5$ m/s (Reynolds $>100{,}000$), proveyendo fricción mecánica para barrer biopelículas en soldaduras orbitales bajo normativa ASME BPE (rugosidad $Ra < 0.4 \mu$m). Las conexiones Tri-Clamp de diafragma se modelan bajo la filosofía SMED (\textit{Single-Minute Exchange of Die}), reduciendo el tiempo de mantenimiento a menos de 3 minutos sin herramientas.

\section{Resultados de las Pruebas de Estrés}

El Gemelo Digital fue sometido a pruebas de estrés variando las condiciones nominales mediante la interfaz Streamlit en \texttt{app.py}.

\subsection{Sobrecarga de Caudal de Alimentación}
Al aumentar la alimentación de soya por encima de 1200 kg/h con un ratio de agua constante (12:1), el flujo hidráulico de lodo a la centrífuga CF-102 supera los 15 m$^3$/h. El motor de restricciones detecta que la fracción de capacidad del decanter excede el límite del 90\%, disparando una alerta roja y paro de emergencia automático (\texttt{EquipmentCapacityError}).

\subsection{Desviación de pH en Lixiviación y Precipitación}
La simulación de una desviación en la dosificación de NaOH en TK-101 que reduce el pH de extracción a 8.0 provocó una disminución de la eficiencia de extracción $\eta_{\text{ext}}$ de 88.0\% a 75.9\%, reduciendo la producción final de polvo ISP a $265.1$ kg/h (una caída del 12.1\% respecto a la línea base de $301.6$ kg/h).

En el precipitador TK-401, una descalibración del sensor de pH que desvíe el pH a 5.2 disminuye la eficiencia de precipitación proteica de 98.0\% a 83.6\%, incrementando la concentración de proteína no precipitada en el suero residual hasta $2.5$ kg/m$^3$.

\subsection{Caída de Temperatura en HTST y Parada HACCP}
Se simuló una pérdida de presión en la línea de vapor de la caldera, disminuyendo la temperatura en el pasteurizador HX-201. Al descender por debajo del límite crítico de control (PCC) de $75^\circ$C, el Gemelo Digital activa un enclavamiento de seguridad que detiene inmediatamente el bombeo hacia el evaporador, redirigiendo el flujo hacia el tanque de reproceso para evitar contaminación microbiológica por \textit{Salmonella}.

\subsection{Pérdida de Vacío en el Evaporador}
Se evaluó el comportamiento del evaporador EV-301 al fallar la bomba de vacío, elevándose la presión interna a $0.70$ bar. El gemelo calcula un aumento en la temperatura de ebullición a $89^\circ$C. Esto activa la penalización cinética por reacción de Maillard, reduciendo el factor de calidad proteica a 0.85 y provocando oscurecimiento estético inaceptable del producto.

\subsection{Beneficio de la Ósmosis Inversa}
La comparación de la simulación térmica pura frente al sistema híbrido (OI + Evaporación) demuestra el beneficio de la pre-concentración por membranas a 24 bar de TMP (Tabla 5).

\begin{table}[H]
\centering
\caption{Comparación Energética: Térmico vs Híbrido (OI)}
\begin{tabular}{lcc}
\toprule
Variable Energética & Sin OI & Con OI \\
\midrule
Agua removida evap. (kg/h) & 8967.6 & 6249.3 \\
Agua removida OI (kg/h) & 0.0 & 2718.3 \\
Potencia Térmica Evap. (kW) & 3530.0 & 2584.5 \\
Ahorro térmico (kW) & 0.0 & 945.5 (25.4\%) \\
Potencia Eléctrica OI (kW) & 0.0 & 9.0 \\
Ahorro Neto Equivalente (kW) & 0.0 & 936.5 \\
\bottomrule
\end{tabular}
\end{table}

El OPEX de la caldera disminuye significativamente, pagando el CAPEX de la unidad de membranas en un periodo inferior a 1.5 años.

\subsection{FMEA Dinámico (D-FMEA)}
El estado del arte de la Industria 5.0 demanda un FMEA Dinámico algorítmico. El Número Prioritario de Riesgo ($NPR = S \times O \times D$) deja de ser un artefacto subjetivo; el Gemelo Digital conecta las APIs telemétricas y los estimadores EKF directamente con las ecuaciones de riesgo:

\begin{itemize}
    \item \textbf{Ocurrencia ($O$) Automatizada:} Incrementos en firmas de vibración o alteraciones en voltaje del variador inducen a los modelos de Deep Learning a estimar desgaste de rodamientos, recalculando continuamente la curva probabilística de Weibull.
    \item \textbf{Severidad ($S$) Condicional:} El DT evalúa la criticidad proyectando la perturbación al futuro. Por ejemplo, si el sello hidráulico del evaporador EV-301 falla y la presión asciende a $0.70$ bar, la temperatura escala de $75^\circ$C a $89^\circ$C, desencadenando cinéticas de Maillard. La severidad se exacerba, inflando el NPR y forzando al MPC a desencadenar protocolos de parada automática.
    \item \textbf{Gestión HACCP:} Ante desviaciones en variables críticas como pérdidas de presión de vapor en los pasteurizadores HTST ($<75^\circ$C), el D-FMEA intercepta el riesgo microbiológico, actuando sobre válvulas de desvío neumáticas en fracciones de segundo.
\end{itemize}

\begin{table}[H]
\centering
\caption{Matriz de Criticidad D-FMEA del Gemelo}
\begin{tabular}{lccccr}
\toprule
Variable de Proceso & S & O & D & NPR & Nivel \\
\midrule
pH Precipitación (TK-401) & 8 & 6 & 3 & 144 & Crítico \\
Humedad Polvo (SD-501) & 9 & 4 & 3 & 108 & Alto \\
Vacío Evaporador (EV-301) & 7 & 5 & 2 & 70 & Alto \\
pH Extracción (TK-101) & 8 & 4 & 2 & 64 & Alto \\
Temp. Pasteur. (HX-201) & 9 & 3 & 1 & 27 & Medio \\
\bottomrule
\end{tabular}
\end{table}

\begin{table}[H]
\centering
\caption{Puntos de Falla y Reacciones Ciberfísicas Autónomas}
\begin{tabular}{lp{2.4cm}p{2.4cm}}
\toprule
Falla & Alerta D-FMEA & Reacción MPC \\
\midrule
Pérdida vacío EV-301 & EKF detecta anomalía, proyecta $89^\circ$C & Corta vapor vivo, bloquea alimentación \\
Colapso secado SD-501 & CFD+ML detecta $\Delta T > 15$ K & Ajuste gas caliente \\
Fouling membranas & LSTM predice $R_f$ creciente & Pre-programa CIP \\
Desvío pH TK-101 & pH baja de 8.75 a 8.0 & Dosificación NaOH correctiva \\
\bottomrule
\end{tabular}
\end{table}

\section{Conclusiones}
El desarrollo e implementación del Gemelo Digital fenomenológico para la planta de Proteína Aislada de Soya (ISP) representa un hito técnico para la transición a la Industria 5.0. La arquitectura propuesta, sustentada por estricta adhesión a las normativas de interoperabilidad AAS (IEC 63278) e ISO 23247, integra modelos fenomenológicos de primer orden con arquitecturas avanzadas de Machine Learning mediante Redes Neuronales Informadas por la Física (PINNs), que fusionan las leyes de conservación termodinámica con la flexibilidad estadística del aprendizaje profundo.

El Filtro de Kalman Extendido (EKF) garantiza la sincronización fidedigna del gemelo digital con la planta física, operando como sensor virtual para inferir variables CTQ no medibles directamente. El motor de Control Predictivo Basado en Modelos (MPC) supera las limitaciones intrínsecas del control PID clásico al gestionar nativamente las dinámicas MIMO acopladas del bioproceso, con visión prospectiva y restricciones mecánicas embebidas en la función de optimización.

Las pruebas de estrés demostraron la capacidad del gemelo para predecir fallas operacionales y actuar proactivamente mediante paradas automáticas por seguridad o inocuidad, mientras el FMEA Dinámico (D-FMEA) recalcula autónomamente el NPR mediante estimaciones de Weibull y simulaciones \textit{What-If}.

El análisis económico y de circularidad integrado valida que la pre-concentración por Ósmosis Inversa, complementada con Ultrafiltración y Diafiltración, ofrece un ahorro neto de 936.5 kW con un \textit{payback} de la inversión inferior a 1.5 años y una mejora en recuperación proteica del 17-26\% frente a la tecnología puramente isoeléctrica, consolidando la viabilidad técnica y ecológica del sistema ciberfísico.

\section*{Referencias}
\begin{enumerate}
    \item Grieves, M., \& Vickers, J. (2017). \textit{Digital Twin: Mitigating Unpredictable, Undesirable Emergent Behavior in Complex Systems}. In Transdisciplinary Perspectives on Complex Systems (pp. 85-113). Springer.
    \item Kritzinger, W., Karner, M., Trautner, G., Jian, J., \& Sihn, W. (2018). \textit{Digital Twin in manufacturing: A categorical literature review and classification}. IFAC-PapersOnLine, 51(11), 1016-1022.
    \item Lusas, E. W., \& Riaz, M. N. (1995). \textit{Soy protein products: processing and use}. The Journal of Nutrition, 125(suppl\_3), 573S-580S.
    \item Kinsella, J. E. (1979). \textit{Functional properties of soy proteins}. Journal of the American Oil Chemists' Society, 56(3), 242-258.
    \item ASME BPE (2019). \textit{Bioprocessing Equipment Standard}. The American Society of Mechanical Engineers.
    \item EHEDG Guidelines (2018). \textit{Hygienic Design Criteria for Sanitary Equipment}. European Hygienic Engineering \& Design Group.
    \item Raissi, M., Perdikaris, P., \& Karniadakis, G. E. (2019). \textit{Physics-informed neural networks: A deep learning framework for solving forward and inverse problems involving nonlinear partial differential equations}. Journal of Computational Physics, 378, 686-707.
    \item Rawlings, J. B., Mayne, D. Q., \& Diehl, M. (2017). \textit{Model Predictive Control: Theory, Computation, and Design}. Nob Hill Publishing.
    \item ISO 23247:2021. \textit{Automation systems and integration --- Digital twin framework for manufacturing}. International Organization for Standardization.
    \item IEC 63278:2023. \textit{Asset Administration Shell for Industrial Applications}. International Electrotechnical Commission.
    \item Mondor, M., Ippersiel, D., Lamarche, F., \& Boye, J. I. (2009). \textit{Effect of electro-acidification treatment and ionic environment on soy protein extract particle size distribution and their influence on membrane fouling when filtered over a wide range of conditions}. Journal of Membrane Science, 231, 209-218.
    \item Alibhai, Z., Mondor, M., Moresoli, C., Ippersiel, D., \& Lamarche, F. (2006). \textit{Production of soy protein concentrates/isolates: traditional and membrane technologies}. Desalination, 191, 351-358.
    \item Roos, Y. H. (1995). \textit{Phase Transitions in Foods}. Academic Press.
    \item Bhandari, B. R., \& Howes, T. (1999). \textit{Implication of glass transition for the drying and stability of dried foods}. Journal of Food Engineering, 40(1-2), 71-79.
    \item Kedem, O., \& Katchalsky, A. (1958). \textit{Thermodynamic analysis of the permeability of biological membranes to non-electrolytes}. Biochimica et Biophysica Acta, 27, 229-246.
    \item OPC Foundation (2022). \textit{OPC UA Companion Specification for the Asset Administration Shell (I4AAS)}. OPC 30270.
    \item Perry, R. H., \& Green, D. W. (2019). \textit{Perry's Chemical Engineers' Handbook}. 9th ed. McGraw-Hill Education.
\end{enumerate}

\end{document}
